\documentclass[12pt]{article}
\usepackage{amsmath, amssymb, amsthm}
\usepackage{hyperref}
\usepackage{geometry}
\geometry{margin=1in}

\title{\textbf{The Gravitational Baseline Principle: Structural Asymmetry and the Pre-Geometric Quantum Substrate}}
\author{
\textbf{Paul Jarvis}\\
Independent Researcher, UK\\
ORCID: 0009-0009-8933-857X\\
\texttt{mrpaulwjarvis@gmail.com}
}
\date{\today}

\begin{document}
\maketitle

\begin{abstract}
We identify a fundamental structural asymmetry between General Relativity (GR) and Quantum Field Theory (QFT): quantum fields require a fixed kinematic background to define locality and operator algebra, whereas GR asserts that the background metric is itself a dynamical physical degree of freedom. We formalize this mismatch as the \emph{Gravitational Baseline Principle}, which states that gravity is not a peer interaction among fields but the universal geometric framework that renders all field-theoretic relations definable. Treating gravity as a conventional quantizable gauge field leads to contradictions in background independence and gauge covariance. To resolve this, we propose a minimal, background-independent quantum substrate from which both smooth manifolds and quantum field algebras co-emerge. Geometry arises from the entanglement structure of a relational density matrix, and the graviton appears as a collective, low-energy excitation of this substrate. We show how the Laplace--Beltrami operator, Einstein dynamics, and thermodynamic gravitational laws emerge in the continuum limit. This framework provides a conceptually coherent path toward unifying GR and QFT without quantizing the metric as a fundamental field.
\end{abstract}

\section{Introduction}
Quantum Field Theory (QFT) describes matter and radiation as operator-valued distributions defined on a fixed background spacetime \cite{Dirac1930}. Locality, microcausality, and the construction of a Fock space all require a predetermined metric structure. General Relativity (GR), by contrast, denies the existence of any prior geometric container: the metric itself is a dynamical field determined by the Einstein equations \cite{Einstein1916}. This architectural mismatch underlies the persistent difficulty in formulating a consistent quantum theory of gravity \cite{DeWitt1967}.

Gravity is universal: all physical systems couple to the metric tensor $g_{\mu\nu}$, which defines causal structure, measurement, and the very notion of locality. Gauge interactions, by contrast, act only on specific internal symmetry representations. This motivates the \emph{Gravitational Baseline Principle}: gravity is the universal geometric framework that renders all other interactions definable. It is not a peer field within a background but the background itself.

In this paper, we develop a minimal pre-geometric quantum substrate from which smooth geometry and quantum fields co-emerge. The substrate is defined by a relational density matrix whose entanglement structure encodes emergent spatial connectivity. We show how the Laplace--Beltrami operator, graviton excitations, and thermodynamic gravitational dynamics arise in the continuum limit.

\section{The Structural Asymmetry Problem}
Einstein's field equations,
\begin{equation}
G_{\mu\nu} = 8\pi G\, T_{\mu\nu},
\end{equation}
imply universal coupling: all fields influence and are influenced by the metric. Gauge interactions, however, operate only on specific internal degrees of freedom.

Perturbative quantization of gravity requires expanding the metric as $g_{\mu\nu} = \eta_{\mu\nu} + h_{\mu\nu}$ and quantizing $h_{\mu\nu}$ as a spin-2 gauge field. But this procedure presupposes a fixed background $\eta_{\mu\nu}$, violating diffeomorphism invariance at the foundational level. This is the core structural asymmetry: QFT requires a background; GR denies one.

\section{The Gravitational Baseline Principle}
Gravity is not a selective interaction but the universal geometric baseline. All physical systems couple non-linearly to the metric, which defines causal structure and measurement. The Planck scales,
\begin{equation}
\ell_P = \sqrt{\frac{\hbar G}{c^3}}, \qquad t_P = \sqrt{\frac{\hbar G}{c^5}},
\end{equation}
mark the breakdown of smooth geometry and the onset of pre-geometric structure.

The Gravitational Baseline Principle states:
\begin{quote}
\emph{Gravity is the universal geometric framework that renders all field-theoretic relations definable. It cannot be quantized as a conventional gauge field without violating background independence.}
\end{quote}

\section{Pre-Geometric Substrate Architecture}
We define the pre-geometric substrate $\mathcal{Q}$ as a background-independent quantum system with state vector $|\Psi\rangle$ in a Hilbert space $\mathcal{H}_{\mathcal{Q}}$. The density matrix $\rho = |\Psi\rangle\langle\Psi|$ encodes relational structure.

Mutual information between subsystems $i$ and $j$,
\begin{equation}
I(i,j) = S(\rho_i) + S(\rho_j) - S(\rho_{ij}),
\end{equation}
defines a symmetric connectivity matrix $A_{ij}$.

\subsection{Emergent Continuum Geometry}
A discrete structural operator acts on test functions $\phi(x)$ mapped to subsystem nodes:
\begin{equation}
\Delta_\mathcal{Q} f_i = \sum_j A_{ij}(f_i - f_j).
\end{equation}
Expanding $\phi(x)$ in a Taylor series and imposing isotropy constraints yields the continuum Laplace--Beltrami operator:
\begin{equation}
\lim_{\alpha\to 0} \Delta_\mathcal{Q} = \nabla^2.
\end{equation}

\subsection{Dynamical Law}
We propose a minimal background-independent evolution equation:
\begin{equation}
\frac{d\rho}{dt} = -i[H_\mathcal{Q}, \rho],
\end{equation}
with Hamiltonian
\begin{equation}
H_\mathcal{Q} = \sum_{i,j} K(A_{ij})\, \sigma_i \otimes \sigma_j,
\end{equation}
where $K$ is a monotonic kernel of entanglement strength. Because $A_{ij}$ depends on $\rho$, the evolution is self-consistent and non-linear.

\section{Emergent Graviton Dynamics}
Small perturbations of the emergent metric,
\begin{equation}
g_{\mu\nu}(x) = \eta_{\mu\nu} + h_{\mu\nu}(x),
\end{equation}
arise from variations in $A_{ij}$ projected onto smooth coordinate functions. After projection onto the transverse-traceless sector, the effective action reduces to the Pauli--Fierz form, demonstrating that the graviton is a collective excitation of the substrate.

\section{Thermodynamic Gravity}
Following Jacobson \cite{Jacobson1995}, Einstein's equations emerge as a thermodynamic equation of state. Fluctuations of the substrate generate higher-curvature corrections:
\begin{equation}
S_{\mathrm{eff}} = \frac{1}{16\pi G} \int d^4x \sqrt{-g}\left(R + \alpha R^2 + \beta R_{\mu\nu}R^{\mu\nu} + \cdots\right),
\end{equation}
with coefficients determined by the entanglement spectrum.

\section{Phenomenological Consequences}
\begin{itemize}
\item \textbf{Holographic noise:} Mutual-information encoding implies an irreducible noise floor in precision interferometry.
\item \textbf{Curvature corrections:} Higher-order terms may affect black hole quasi-normal modes and early-universe cosmology.
\item \textbf{Modified dispersion:} Collective excitations satisfy $\omega^2 = k^2[1 + \epsilon(k)]$, with $\epsilon(k)$ suppressed by the microscopic scale.
\end{itemize}

\section{Relation to Existing Approaches}
Unlike Loop Quantum Gravity or Causal Set Theory, which quantize classical geometric structures, this framework treats geometry as emergent from non-spatial quantum information. It shares conceptual ground with ER=EPR \cite{Susskind2016} and entanglement-induced geometry \cite{VanRaamsdonk2010}, but provides a concrete dynamical substrate.

\section{Limitations and Outlook}
This framework is conceptual and minimalistic. A full treatment requires:
\begin{itemize}
\item explicit reconstruction of Lorentzian signature,
\item derivation of matter field algebras,
\item analysis of renormalization behavior,
\item numerical simulations of substrate dynamics.
\end{itemize}
Nonetheless, the model provides a coherent path toward unifying GR and QFT without quantizing the metric.

\section{Conclusion}
The Gravitational Baseline Principle resolves the structural asymmetry between GR and QFT by recognizing gravity as the universal geometric framework rather than a conventional gauge field. A pre-geometric quantum substrate with relational entanglement structure yields emergent geometry, graviton excitations, and thermodynamic gravitational dynamics. This approach offers a conceptually consistent foundation for quantum gravity.

\begin{thebibliography}{9}

\bibitem{DeWitt1967}
B.~S.~DeWitt, ``Quantum Theory of Gravity,'' \emph{Phys. Rev.} \textbf{160}, 1113 (1967).

\bibitem{Dirac1930}
P.~A.~M.~Dirac, \emph{The Principles of Quantum Mechanics}, Oxford University Press (1930).

\bibitem{Einstein1916}
A.~Einstein, ``The Foundation of the General Theory of Relativity,'' \emph{Annalen der Physik} \textbf{49}, 769 (1916).

\bibitem{Jacobson1995}
T.~Jacobson, ``Thermodynamics of Spacetime: The Einstein Equation of State,'' \emph{Phys. Rev. Lett.} \textbf{75}, 1260 (1995).

\bibitem{Susskind2016}
L.~Susskind, ``ER=EPR, GHZ, and the Consistency of Quantum Measurements,'' \emph{Fortschritte der Physik} \textbf{64}, 72 (2016).

\bibitem{VanRaamsdonk2010}
M.~Van Raamsdonk, ``Building up spacetime with quantum entanglement,'' \emph{Gen. Rel. Grav.} \textbf{42}, 2323 (2010).

\end{thebibliography}

\end{document}
